\section{Results II: post-release data}
\label{sec:postrelease}

\begin{figure}[!htbp]
\centering
\pgfplotstableread{
x y em ep fam res
1.1096 0 0.0608 0.0688 1 1
1.0004 1 0.0006 0.0006 1 1
0.9979 2 0.0013 0.0013 1 1
0.4458 3 0.0143 0.0145 1 1
0.6159 4 0.0247 0.0266 1 1
0.9799 5 0.0438 0.0408 1 0
0.9894 6 0.0133 0.0125 0 0
0.9667 7 0.0266 0.0300 0 0
0.9509 8 0.0334 0.0372 0 0
0.9528 9 0.0290 0.0280 0 1
0.8358 10 0.0502 0.0462 0 1
}\forestdata
\pgfplotsset{forest ticks/.style={ytick={0,1,2,3,4,5,6,7,8,9,10}, yticklabels={{H5: TimesFM-3 uni vs Naive, FX and FRED},{H4d: multivariate vs univariate, spurious covariate},{H4c: multivariate vs univariate, independent series},{H4b: multivariate vs univariate, known covariate},{H4a: multivariate vs univariate, lead-lag 24},{H2: W-seen vs W-post DiD, TimesFM-3 vs Chronos-2},{P5: TimesFM-3 eval vs TimesFM-3},{P4a: TimesFM-3 vs TimesFM-3 uni},{P3: TimesFM-3 uni vs TimesFM-2.5},{P2: TimesFM-3 uni vs Chronos-2 uni},{P1: TimesFM-3 uni vs Stat.\ ensemble}}, ymin=-0.7, ymax=10.7}}

\begin{tikzpicture}
\begin{axis}[
  width=0.62\textwidth, height=9.2cm,
  xmode=log, log basis x=10, xmin=0.4, xmax=1.3,
  xtick={0.5,0.6,0.7,0.8,0.9,1,1.1,1.2}, xticklabels={0.5,0.6,0.7,0.8,0.9,1,1.1,1.2},
  forest ticks, yticklabel style={font=\scriptsize, align=right},
  xlabel={ratio of errors, A / B (log scale; $<1$: A better)}, xlabel style={font=\footnotesize},
  tick label style={font=\footnotesize},
  axis lines*=left, ymajorgrids=false, xmajorgrids=true, grid style={black!10},
  clip=false]
\addplot[black!40, dashed] coordinates {(1,-0.7) (1,10.7)};
\addplot[only marks, mark=*, mark size=1.6pt, black,
  error bars/.cd, x dir=both, x explicit, error bar style={line width=0.8pt}]
  table[x=x, y=y, x error minus=em, x error plus=ep] {\forestdata};
\end{axis}
\end{tikzpicture}
\caption{Primary (family P, MASE) and hypothesis (family H) contrasts: estimate and two-way cluster bootstrap 95\% interval
(unadjusted). P-tests are on post-release Tier B data: P1--P3 without covariates on the same series and windows (context lengths differ on
hourly data), P4a on the group tasks (no covariates), P5 between two configurations with the covariates of load and river flow; H4a--H4d are on controlled
multivariate data; H5 is on exchange rates, Treasury yields and oil prices. H3 is omitted (its cell never activates the mechanism
it tests, Section~\ref{sec:controlled}). Decisions use Holm-adjusted $p$-values (Tables~\ref{tab:familyp} and~\ref{tab:familyh}).}
\label{fig:forest}
\end{figure}


\paragraph{Primary contrasts (family P).} On the four structured domains of Tier B, with no covariates for any model, TimesFM~3 in
univariate mode has a MASE \PoneMase{} times that of the statistical ensemble (95\% CI \PoneMaseLo--\PoneMaseHi) and an SQL
\PoneSql{} times (\PoneSqlLo--\PoneSqlHi); both differences are resolved (Table~\ref{tab:familyp}, Figure~\ref{fig:forest}). Against
the other foundation models the margins are an order of magnitude smaller: \PtwoMase{} (MASE) and \PtwoSql{} (SQL) against
Chronos-2, both resolved, and \PthreeSql{} (SQL, resolved) and \PthreeMase{} (MASE, not resolved after Holm correction) against its
predecessor TimesFM~2.5. Joint forecasting of the real group tasks (P4a: price, FX, FRED) lowers SQL by about 3\% (resolved; the
MASE difference is not resolved). The flags of Google's benchmark runners (P5) change neither metric by more than about 1\%, at
twice the compute; neither difference is resolved.

\begin{table}[!htbp]
\centering
\caption{Family P (primary): information-matched contrasts on Tier B. Ratio of errors, A / B ($<1$: A better), with the
two-way cluster bootstrap interval and the series-only interval; decisions by Holm-adjusted $p$ over the ten tests.}
\label{tab:familyp}
\scriptsize
\setlength{\tabcolsep}{4pt}
\begin{tabular}{lllrllll}
\toprule
Test & Contrast (A vs B) & Metric & Ratio & 95\% CI (two-way) & 95\% CI (series) & $p_{\mathrm{Holm}}$ & Decision \\
\midrule
P1 & TimesFM-3 uni vs Stat.\ ensemble & MASE & 0.836 & [0.786, 0.882] & [0.804, 0.867] & 0.001 & resolved \\
P1 & TimesFM-3 uni vs Stat.\ ensemble & SQL & 0.707 & [0.649, 0.764] & [0.655, 0.757] & 0.001 & resolved \\
P2 & TimesFM-3 uni vs Chronos-2 uni & MASE & 0.953 & [0.924, 0.981] & [0.937, 0.969] & 0.014 & resolved \\
P2 & TimesFM-3 uni vs Chronos-2 uni & SQL & 0.956 & [0.933, 0.979] & [0.942, 0.971] & 0.003 & resolved \\
P3 & TimesFM-3 uni vs TimesFM-2.5 & MASE & 0.951 & [0.917, 0.988] & [0.924, 0.981] & 0.066 & not resolved \\
P3 & TimesFM-3 uni vs TimesFM-2.5 & SQL & 0.939 & [0.914, 0.965] & [0.916, 0.963] & 0.001 & resolved \\
P4a & TimesFM-3 vs TimesFM-3 uni & MASE & 0.967 & [0.940, 0.997] & [0.947, 0.986] & 0.105 & not resolved \\
P4a & TimesFM-3 vs TimesFM-3 uni & SQL & 0.966 & [0.947, 0.988] & [0.950, 0.982] & 0.025 & resolved \\
P5 & TimesFM-3 eval vs TimesFM-3 & MASE & 0.989 & [0.976, 1.002] & [0.983, 0.995] & 0.106 & not resolved \\
P5 & TimesFM-3 eval vs TimesFM-3 & SQL & 0.991 & [0.982, 1.000] & [0.986, 0.996] & 0.106 & not resolved \\
\bottomrule
\end{tabular}

\end{table}

\paragraph{By domain (S1).} The advantage over the statistical ensemble is resolved in every domain, from 11\% lower MASE on
electricity load to 22\% on river flow (Table~\ref{tab:secondary}). Against Chronos-2 the per-domain differences are resolved only
on weather; against TimesFM~2.5 on load and weather. The cross-model margins are therefore not uniform across domains, and the
pooled P2 result rests mostly on the largest domain (weather: 2{,}280 of 3{,}855 series-windows).

\paragraph{H1: is the advantage smaller after the release?} With the same configuration as on Tier A (package defaults, covariates
included), TimesFM~3's MASE relative to our statistical ensemble on Tier B is \HoneBMase{} (\HoneBMaseLo--\HoneBMaseHi) and its SQL
\HoneBSql{} (\HoneBSqlLo--\HoneBSqlHi), against \HoneAMase{} and \HoneASql{} relative to fev's ensemble on the benchmark tasks.
Descriptively, the point-forecast advantage is about a quarter smaller on post-release data (18\% against 25\%) and the
probabilistic advantage almost unchanged (31\% against 34\%). The two tiers differ in datasets, horizons and ensemble members, so
this comparison carries no test and cannot separate exposure from difficulty; it says that the benchmark-level advantage over
classical methods survives on new data, slightly reduced for the median forecast.

\paragraph{H5: near-random-walk series.} On exchange rates, Treasury yields and oil prices, TimesFM~3's MASE is \Hfive{} times that
of the random walk (\HfiveLo--\HfiveHi): H5 is supported, and more strongly than stated, since the model is \emph{worse} than the
naive forecast rather than not much better. It is also worse than the statistical ensemble on these series (\NrwStatEns,
\NrwStatEnsLo--\NrwStatEnsHi; descriptive), while the difference from Chronos-2 is not resolved (\NrwChronos,
\NrwChronosLo--\NrwChronosHi). A random walk is the optimal forecast for a martingale, so no model can beat it by a resolved margin
without an evaluation problem; the negative controls in Tier C behave the same way (Section~\ref{sec:controlled}).

\paragraph{Covariates and context length (S3, S4).} Supplying the observed future temperature for load and the observed future
precipitation for river flow, an oracle that no real forecast has, does not resolve a gain: \texttt{tfm3} against \texttt{tfm3-uni}
gives 0.974 (0.904--1.060), and Chronos-2 with and without covariates 0.969 (0.930--1.008). On hourly data, cutting TimesFM~3's
context to the 2{,}016 points the classical models use raises its MASE by 4\% (S4, resolved); the long context helps, but explains
only a small part of the gap to the classical ensemble on hourly series (11--13\% in load and price).

\paragraph{Memorisation probe (H2, S2).} TimesFM~3's MASE relative to Chronos-2 is 0.939 on the 2023 W-seen windows and 0.959 on the
same pages in 2026; the difference in differences is \Htwo{} (\HtwoLo--\HtwoHi), not resolved (Table~\ref{tab:familyh}). The design
had 80\% power to detect a difference of about 5\%, so a memorisation advantage of that size is unlikely; smaller effects cannot be
excluded, and a null difference does not show that the pages were not memorised. The secondary contrasts agree
(Table~\ref{tab:secondary}): with the statistical ensemble or TiRex-2 as the reference, and for TimesFM~2.5, the difference is not
resolved. On pages born after 2023 (W-new), TimesFM~3's advantage over Chronos-2 is smaller than on W-post (0.984 against 0.959); the
raw contrast is not resolved, and the contrast adjusted for context features (0.942, 0.897--0.996) is not resolved after Holm
correction. If anything, TimesFM~3 does slightly better on long-lived pages than on new ones, which is consistent with familiarity
with the type of series rather than with the particular windows.

\paragraph{Calibration (RQ4).} The empirical coverage of the nominal 80\% interval on Tier B is close to nominal for TimesFM~3
(\CovBtfm{}), TimesFM~2.5 and TiRex-2, lower for Chronos-2 (0.750) and lower for the statistical ensemble (\CovBstat{}), whose MSTL
intervals ignore the uncertainty of the seasonal component (Table~\ref{tab:calibration}). On Wikipedia pageviews all foundation
models are too narrow (coverage 0.72--0.76) and the ensemble far too wide (0.95). The Wilson intervals use the number of series as the sample size and are therefore wide; they are a guide,
not a test.

\paragraph{Exploratory: spikes and lineage.} These analyses were declared exploratory (protocol §5.3) and are not corrected for
multiplicity. (i) \emph{Spike anticipation on W.} In 191 of the 5{,}588 Wikipedia windows the observed maximum exceeds three times
the 90th percentile of the context. In the 161 of them where the spike starts after the forecast origin, no foundation model
anticipates it: the maximum of TimesFM~3's 0.9 quantile over the window has a median of 1.09 times the context percentile (observed:
4.7), and exceeds three times the percentile in 9\% of these windows (Chronos-2 8\%, TiRex-2 and TimesFM~2.5 3\%). The naive
forecast reaches that level more often (48\%) only because its intervals widen with the horizon. In the 30 windows where the spike is
already under way at the origin, all foundation models extend it. The share does not differ between W-seen and W-post for TimesFM~3
(0.189 and 0.191), which agrees with the memorisation probe. (ii) \emph{Lineage.} On the structured domains of Tier B,
TimesFM~3's MASE relative to the statistical ensemble is 0.857 (0.823--0.889) on series whose source has earlier-period data in
public pre-training corpora (load, price, weather) and 0.776 (0.636--0.922) on rivers, for which none is known. The advantage is not
larger where exposure to the source is documented; since the lineage groups coincide with domains, this is not a test of exposure.

\begin{table}[!htbp]
\centering
\caption{Calibration: mean empirical coverage of the central 80\% and 60\% intervals (protocol §5.6). The Wilson interval treats
each series as one observation (each series' coverage averaged over its windows and steps).}
\label{tab:calibration}
\small
\begin{tabular}{lrrr}
\toprule
Model & Series & Coverage 80\% & Coverage 60\% \\
\midrule
\multicolumn{4}{l}{\textit{Tier B (structured, univariate tasks)}} \\
TimesFM-3 uni & 84 & 0.796 & 0.596 \\
TimesFM-3 & 84 & 0.802 & 0.598 \\
TimesFM-3 eval & 84 & 0.807 & 0.604 \\
TimesFM-2.5 & 84 & 0.807 & 0.607 \\
Chronos-2 uni & 84 & 0.750 & 0.546 \\
TiRex-2 uni & 84 & 0.797 & 0.599 \\
Stat.\ ensemble & 84 & 0.779 & 0.630 \\
\multicolumn{4}{l}{\textit{Tier W, W-post}} \\
TimesFM-3 & 190 & 0.761 & 0.553 \\
TimesFM-2.5 & 190 & 0.755 & 0.563 \\
Chronos-2 & 190 & 0.722 & 0.520 \\
TiRex-2 & 190 & 0.760 & 0.560 \\
Stat.\ ensemble & 190 & 0.946 & 0.886 \\
\multicolumn{4}{l}{\textit{Tier C (all cells)}} \\
TimesFM-3 & 7600 & 0.807 & 0.597 \\
TimesFM-2.5 & 7600 & 0.804 & 0.616 \\
Chronos-2 & 7600 & 0.780 & 0.544 \\
TiRex-2 & 7200 & 0.821 & 0.635 \\
Stat.\ ensemble & 7600 & 0.866 & 0.728 \\
\multicolumn{4}{l}{\textit{Tier P, both issues: load, price, river}} \\
TimesFM-3 uni & 44 & 0.699 & --- \\
TimesFM-2.5 & 44 & 0.707 & --- \\
Chronos-2 uni & 44 & 0.672 & --- \\
TiRex-2 uni & 44 & 0.732 & --- \\
Stat.\ ensemble & 44 & 0.876 & --- \\
\multicolumn{4}{l}{\textit{Tier P, both issues: Wikipedia}} \\
TimesFM-3 uni & 318 & 0.687 & --- \\
TimesFM-2.5 & 318 & 0.694 & --- \\
Chronos-2 uni & 318 & 0.670 & --- \\
TiRex-2 uni & 318 & 0.703 & --- \\
Stat.\ ensemble & 318 & 0.944 & --- \\
\bottomrule
\end{tabular}

\end{table}

\begin{table}[!htbp]
\centering
\caption{Secondary families (Holm within each family). S1: family P contrasts by domain; S2: Wikipedia memorisation probe; S3:
oracle future covariates; S4: TimesFM~3 restricted to the classical models' context.}
\label{tab:secondary}
\scriptsize
\begin{tabular}{llrlll}
\toprule
Test & Contrast (A vs B) & Ratio & 95\% CI & $p_{\mathrm{Holm}}$ & Resolved \\
\midrule
\multicolumn{6}{l}{\textit{S1: family P contrasts by domain (MASE)}} \\
P1, load & TimesFM-3 uni vs Stat.\ ensemble & 0.892 & [0.824, 0.950] & 0.004 & yes \\
P2, load & TimesFM-3 uni vs Chronos-2 uni & 0.960 & [0.909, 1.005] & 0.245 & no \\
P3, load & TimesFM-3 uni vs TimesFM-2.5 & 0.925 & [0.882, 0.965] & 0.008 & yes \\
P1, price & TimesFM-3 uni vs Stat.\ ensemble & 0.872 & [0.781, 0.952] & 0.008 & yes \\
P2, price & TimesFM-3 uni vs Chronos-2 uni & 0.979 & [0.930, 1.029] & 0.817 & no \\
P3, price & TimesFM-3 uni vs TimesFM-2.5 & 0.963 & [0.933, 0.992] & 0.069 & no \\
P1, weather & TimesFM-3 uni vs Stat.\ ensemble & 0.808 & [0.763, 0.850] & 0.001 & yes \\
P2, weather & TimesFM-3 uni vs Chronos-2 uni & 0.956 & [0.929, 0.983] & 0.026 & yes \\
P3, weather & TimesFM-3 uni vs TimesFM-2.5 & 0.961 & [0.939, 0.982] & 0.011 & yes \\
P1, river & TimesFM-3 uni vs Stat.\ ensemble & 0.776 & [0.628, 0.923] & 0.015 & yes \\
P2, river & TimesFM-3 uni vs Chronos-2 uni & 0.917 & [0.834, 1.000] & 0.201 & no \\
P3, river & TimesFM-3 uni vs TimesFM-2.5 & 0.955 & [0.835, 1.107] & 0.817 & no \\
\midrule
\multicolumn{6}{l}{\textit{S2: Wikipedia (MASE; DiD: W-seen minus W-post advantage)}} \\
DiD, reference stat.\ ensemble & TimesFM-3 vs Stat.\ ensemble & 1.009 & [0.915, 1.108] & 1.000 & no \\
DiD, reference TiRex-2 & TimesFM-3 vs TiRex-2 & 0.986 & [0.951, 1.024] & 1.000 & no \\
DiD for TimesFM-2.5 & TimesFM-2.5 vs Chronos-2 & 0.980 & [0.936, 1.021] & 1.000 & no \\
W-post vs W-new & TimesFM-3 vs Chronos-2 & 0.974 & [0.941, 1.009] & 0.737 & no \\
adjusted DiD & TimesFM-3 vs Chronos-2 & 1.013 & [0.971, 1.056] & 1.000 & no \\
adjusted W-post vs W-new & TimesFM-3 vs Chronos-2 & 0.942 & [0.897, 0.996] & 0.220 & no \\
\midrule
\multicolumn{6}{l}{\textit{S3: oracle covariates, load and river (oracle information condition); S4: classical context}} \\
P4b & TimesFM-3 vs TimesFM-3 uni & 0.973 & [0.904, 1.060] & 0.467 & no \\
Chronos-2 covariates & Chronos-2 vs Chronos-2 uni & 0.969 & [0.930, 1.008] & 0.217 & no \\
S4, hourly domains & TimesFM-3 c2016 vs TimesFM-3 uni & 1.040 & [1.011, 1.072] & 0.004 & yes \\
\bottomrule
\end{tabular}

\end{table}
